\section{Bussen och routingen}
\label{sec:bus}

Två noder kan kopplas direkt till varandra. Fem kan inte --- det blir tio förbindelser, och antalet
växer kvadratiskt. Lösningen är en \textbf{delad buss}: alla noder på samma ledning.

Det får en omedelbar konsekvens: \textbf{alla hör allt}. En buss är broadcast, och det är inget man
väljer bort. Vill man att ett meddelande ska nå en viss nod får man skriva det i meddelandet.

\subsection{Var routingen ligger}

Bussens enda jobb är att leverera varje byte till varje nod. Den läser inte adresser, den känner
inte till frames, den vet inte ens att det finns ett protokoll.

\textbf{Routingen ligger i noden.} Varje nod tar emot varje byte, matar den till sin parser, och
när en komplett frame extraherats jämför den framens \code{DST} med sin egen adress:

\begin{cppcode}
void Node::tick() noexcept
{
    if (!myParser.isFrameReady()) { return; }

    Frame frame{};
    if (!myParser.extractFrame(frame)) { myParser.reset(); return; }
    myParser.reset();

    if (frame.dstAddr != myAddr) { return; }   // not for us
    handleFrame(frame);
}
\end{cppcode}

Att lägga routingen i noden i stället för i bussen är inte en förenkling för kursens skull. Det är
hur en riktig buss fungerar: ledningen kan inte läsa adresser.

\subsection{DST och SRC}

\code{DST} säger vem framen är till. \code{SRC} säger vem den kom från, och behövs för att
mottagaren ska kunna svara --- utan den vet den inte vart svaret ska.

\code{SEQ} parar ihop svar med fråga. En nod som skickat tre förfrågningar och får ett svar måste
kunna avgöra vilken av dem det svarar på.

\begin{aside}
\note CAN gör det här helt annorlunda: där finns ingen destinationsadress alls, och
identifieraren beskriver meddelandets \emph{innehåll} i stället för dess mottagare. Jämförelsen
görs i \kapref{ch:can}, och den är en av kursens mest lärorika.
\end{aside}
